\chapter{ThermoDynamics of Blackbody Radiation}\label{ch:blackbody}

\small
\begin{quote}\itshape
A perfectly stable cavity, at rest at temperature $T$ in empty space,
fills with electromagnetic radiation of a spectrum that depends only on
$T$. The mystery is not the spectrum --- Planck wrote it down in 1900.
The mystery is the cutoff: why the high-frequency tail does not
diverge.
\end{quote}
\begin{quote}\itshape
The whole of the energy of the radiation is degraded into heat at a
rate that is finite, never infinite. (Planck, paraphrased)
\end{quote}
\normalsize

\section{The Mystery of the Cutoff}

The classical electromagnetic theory of a cavity in equilibrium with
its walls assigns equal energy $k_{\mathrm{B}}T$ to every standing-wave
mode. The number of modes between frequencies $\omega$ and $\omega +
d\omega$ grows as $\omega^2 d\omega$, so the predicted spectral energy
density grows as $\omega^2 k_{\mathrm{B}}T$ and the total energy
$\int_0^\infty \omega^2 k_{\mathrm{B}}T\,d\omega$ diverges --- the
\emph{ultraviolet catastrophe}. Real cavities of course hold finite
energy. The puzzle, in 1900, was where the high-frequency end of the
spectrum gets cut off.

Planck's resolution introduced an inseparable quantum of energy
$\hbar\omega$ per mode, with the Boltzmann factor
$\exp(-\hbar\omega/k_{\mathrm{B}}T)$ suppressing modes for which
$\hbar\omega \gg k_{\mathrm{B}}T$. The spectrum becomes
\[
  u(\omega,T) \;=\; \frac{\hbar\omega^3/\pi^2 c^3}{\exp(\hbar\omega/k_{\mathrm{B}}T) - 1},
\]
finite for every $T$, with a sharp drop at $\omega \sim
k_{\mathrm{B}}T/\hbar$ that places the cutoff in the right physical
place. The mathematics works. The physical picture --- \emph{what
happens to the high-frequency radiation that would have been there
without the cutoff} --- was less clear at the time and is still treated
in many textbooks as merely a statistical consequence of mode
counting.

\section{The Cutoff as Dissipation}

Read in the with-Dynamics framing of this volume, the Planck cutoff is
the radiative analogue of the viscous cutoff in turbulent flow.

In turbulent compressible flow at high Reynolds number, the
deterministic continuum equations would describe a flow with structure
on all scales down to zero. The flow that is actually realised, in
finite resolution, has a cutoff at the viscous scale --- the scale
below which the molecular viscosity becomes comparable to the inertial
forces and the energy that would have cascaded to still smaller scales
is instead converted to internal energy. The dissipation
$D = \int\!\int \mu_{\mathrm{mom}}|\nabla u|^2\,dx\,dt$ of this volume
is the cumulative rate of that conversion. Without the cutoff the
cascade would run to infinity; with the cutoff it terminates by
turning small-scale kinetic energy into heat.

The same pattern governs the radiation field in a cavity. The
classical (unquantised) Maxwell theory predicts modes at arbitrarily
high frequency, each carrying its equipartition $k_{\mathrm{B}}T$. The
field that is actually realised has a cutoff at the quantum scale ---
the scale at which a single quantum $\hbar\omega$ exceeds the thermal
energy $k_{\mathrm{B}}T$ and the corresponding modes can no longer be
populated. The high-frequency radiation that would have been there
without the cutoff is absorbed by the walls and converted to heat ---
exactly the radiative counterpart of the small-scale turbulent kinetic
energy being absorbed at the viscous scale and converted to heat.

\begin{center}
\begin{tabular}{l|l}
\textbf{Turbulence (RNS)} & \textbf{Blackbody radiation} \\
\hline
inertial flow at small scales & high-frequency modes \\
viscous cutoff at $\ell \sim (\nu^3/\epsilon)^{1/4}$ &
quantum cutoff at $\omega \sim k_{\mathrm{B}}T/\hbar$ \\
small-scale kinetic energy $\to$ heat &
above-cutoff radiation $\to$ heat \\
dissipation $D = \int\mu|\nabla u|^2\,dx\,dt \geq 0$ &
absorption rate $\geq 0$ \\
2nd Law as inevitability of $D > 0$ &
2nd Law as inevitability of absorption \\
\end{tabular}
\end{center}

The cutoff is what makes both phenomena \emph{finite}. Without it, the
cascade would run unbounded and the field would carry infinite energy.
The mechanism is the same in spirit: a structural barrier below which
the small-scale degrees of freedom can no longer be supported, and
above which the energy that would have populated them is delivered
instead to the thermal bath.

\section{Computational Blackbody}

The same computational framing that organises this volume applies. A
finite-precision solver of the Maxwell equations on a cavity at
temperature $T$ has a natural cutoff set by its spatial resolution: it
cannot represent modes with wavelength shorter than its mesh spacing,
and the radiation energy that would have lived in those modes is
re-deposited, by the discretisation, as a heating source in the bath.
The Planck cutoff is, in this reading, not a quantum-mechanical
postulate but the cutoff that any finite-resolution representation of
the field must impose --- realised in Nature by whatever physical
process limits the resolution at very short wavelength, and realised
on a computer by the mesh.

This view is developed in detail in the author's earlier
\emph{Computational Black-Body Radiation}~\cite{blackbody}. The
present chapter records that the Planck cutoff and the viscous cutoff
of turbulent RNS are two instances of the same with-Dynamics
mechanism: small-scale degrees of freedom that the underlying
continuum equations \emph{would} support, and that finite resolution
\emph{does not}, deliver their energy as heat. The 2nd Law, in this
joint reading, is the inevitability of that delivery --- in
turbulence, in radiation, and in every other context in which a
continuum field is realised at finite resolution.
